\begin{exercise}[A conserved current from a Killing field]{Tong, Cambridge Part III Sheet 3}
Let \(T^{\mu\nu}\) be symmetric and covariantly conserved, and let
\(\xi^\mu\) be Killing.  Prove that \(J^\mu=T^{\mu\nu}\xi_\nu\) is conserved.
Explain when the associated hypersurface integral is independent of the
chosen hypersurface.
\end{exercise}

\begin{exercise}[Electromagnetic stress and the Lorentz force]{MIT 8.962, Problem Set 5}
Starting from the Maxwell stress tensor in curved spacetime, use Maxwell's
equations to show
\[
\nabla_\mu T_{\rm EM}^{\mu\nu}=-F^\nu{}_\rho J^\rho.
\]
Interpret the result together with conservation of total stress energy.
\end{exercise}

\begin{exercise}[Null dust follows null geodesics]{Cambridge Part III, General Relativity examples}
For \(T^{\mu\nu}=\rho k^\mu k^\nu\), with \(k^\mu k_\mu=0\), use
\(\nabla_\mu T^{\mu\nu}=0\) to derive both the propagation of \(\rho\) and the
geodesic equation for \(k^\mu\), allowing for a nonaffine parameter.
\end{exercise}

\begin{exercise}[Sound speed and relativistic causality]{Cambridge Part III, General Relativity examples}
Linearize number conservation and \(\nabla_\mu T^{\mu\nu}=0\) for a
barotropic perfect fluid about a uniform rest state in Minkowski spacetime.
Derive the sound-wave equation and identify the condition imposed by
causality.
\end{exercise}

\begin{exercise}[Relativistic Bernoulli constant]{McGreevy, Problem Set 6}
For an isentropic stationary perfect fluid with specific enthalpy
\(h=(\rho+p)/n\) in a spacetime with timelike Killing field \(\xi^\mu\), show
that \(h\,u_\mu\xi^\mu\) is constant along each streamline.
\end{exercise}

\begin{exercise}[Liouville's theorem and surface brightness]{MIT 8.962, Problem Set 11}
Use the collisionless Boltzmann equation on the null mass shell to prove that
photon occupation number is constant along a ray.  Deduce that
\(I_\nu/\nu^3\) is conserved and derive the bolometric surface-brightness
dimming between emitter and observer.
\end{exercise}

\begin{exercise}[Transport of polarization]{Cambridge Part III, 2018 exam}
Carry the geometric-optics expansion of Maxwell theory one order beyond the
eikonal equation.  Show that the polarization vector is parallel transported
along the null rays, up to the gauge freedom of adding a multiple of the ray
tangent.
\end{exercise}

\begin{exercise}[Energy conditions for a scalar field]{Tong, Cambridge Part III Sheet 3}
For a minimally coupled real scalar with potential \(V(\phi)\), evaluate
\(T_{\mu\nu}v^\mu v^\nu\) for timelike \(v^\mu\).  Determine when the weak,
null, and strong energy conditions hold, and explain how positive potential
energy can violate the strong condition.
\end{exercise}

\begin{exercise}[Stress tensor of a point particle]{MIT 8.962, Problem Set 9}
Vary the worldline action \(I=-m\int\dd\tau\) with respect to the metric and
derive the distributional stress tensor of a point particle.  Show that its
covariant conservation is equivalent to the particle's geodesic equation.
\end{exercise}

\begin{exercise}[Why Maxwell theory is special in four dimensions]{Cambridge Part III, 2025 exam}
Compute the trace of the Maxwell stress tensor in \(d\) spacetime dimensions.
Relate its vanishing to conformal invariance of the source-free Maxwell
action, and show why \(d=4\) is distinguished.
\end{exercise}
